\chapter{Guide de reproduction et éléments techniques}
\section{Préparation de l'environnement}
Le dépôt utilise Python 3.12 ou 3.13. L'installation de base contient les dépendances de parsing et d'export. Les options vectorielles, Docling et extraction peuvent être ajoutées selon le scénario. Les commandes suivantes sont fournies à titre de reproduction technique; elles doivent être exécutées dans un environnement autorisé et avec les configurations d'accès appropriées.

\begin{lstlisting}[caption={Création d'un environnement et installation des dépendances}]
python -m venv .venv
.\.venv\Scripts\activate
pip install -e ".[dev]"
pip install -e ".[dev,vector]"
pip install -e ".[docling]"
\end{lstlisting}

\section{Exécution de l'évaluation}
L'évaluateur compare les sorties à des gabarits de référence. Les seuils peuvent être utilisés pour faire échouer une exécution d'intégration continue lorsque le niveau de qualité minimal n'est pas atteint.

\begin{lstlisting}[caption={Exemple de commande d'évaluation}]
python scripts\run_evaluation.py \
  --k 1 3 5 \
  --min-f1 0.90 \
  --min-value-accuracy 0.95
\end{lstlisting}

\section{Exemple anonymisé de flux de données}
La représentation ci-dessous illustre la nature des objets produits. Les valeurs sont fictives et ne correspondent à aucune donnée de patient.

\begin{lstlisting}[caption={Exemple de candidat exportable anonymisé}]
{
  "canonical_key": "creatinine",
  "value": "61.9",
  "unit": "umol/L",
  "target_column": "CREATININE",
  "confidence": 0.98,
  "evidence": "Creatinine .... 61.9 umol/L",
  "status": "validated"
}
\end{lstlisting}

\section{Politique d'écriture}
\begin{longtable}{p{0.23\textwidth}p{0.67\textwidth}}
\toprule
Politique & Comportement\\
\midrule
\texttt{empty\_only} & Écrit une proposition uniquement si la cellule cible est vide. Cette politique est la plus prudente pour une première intégration.\\
\texttt{always} & Autorise la mise à jour d'une cellule existante lorsque l'exécution est explicitement configurée ainsi.\\
\texttt{never} & N'écrit aucune valeur dans une cellule déjà renseignée.\\
\bottomrule
\end{longtable}

\section{Pistes de backlog}
Les éléments suivants ont été identifiés comme évolutions possibles :
\begin{itemize}
  \item exposer une API HTTP pour déclencher la pipeline et gérer des traitements asynchrones ;
  \item pseudonymiser les informations identifiantes avant toute indexation ;
  \item intégrer une boucle de revue humaine avec statuts brouillon, confirmé et rejeté ;
  \item enrichir les mesures de qualité par des métriques de latence et de supervision ;
  \item étendre le jeu de test à de nouvelles structures de documents et à un protocole de validation externe.
\end{itemize}

\chapter{Tableaux complémentaires}
\section{Interprétation des indicateurs}
\begin{longtable}{p{0.25\textwidth}p{0.65\textwidth}}
\toprule
Indicateur & Interprétation dans ce rapport\\
\midrule
Précision & Part des champs proposés qui sont effectivement attendus par le gabarit de référence.\\
Rappel & Part des champs attendus qui ont été proposés par le système.\\
F1 & Moyenne harmonique de la précision et du rappel; elle décrit la présence des champs, non la justesse de chaque valeur.\\
Exactitude des valeurs & Part des valeurs proposées qui correspondent aux valeurs de référence après normalisation.\\
Faux autoremplissage & Nombre de champs remplis alors que le gabarit de référence les laisse vides.\\
Taux de validation & Part des propositions qui respectent les règles de type, de chemin et d'unité déclarées.\\
\bottomrule
\end{longtable}

\section{Périmètre de confidentialité}
Le rapport a été rédigé de manière à ne contenir aucune donnée clinique identifiable, aucune capture d'écran d'outil interne et aucune affirmation sur l'activité d'AMALYTICS qui ne serait pas fournie dans le contexte du stage. Les figures sont des diagrammes conceptuels créés pour ce mémoire. Les exemples d'export sont fictifs.
